\subsection{Related Work}
Our work sits at the intersection of three distinct streams of literature: stochastic inventory control (specifically Deep Reinforcement Learning approaches), the application of Large Language Models (LLMs) to optimization, and the evaluation of interactive dialogue systems through user simulation.

\subsubsection{Stochastic Inventory Control and Deep Reinforcement Learning.}
The theoretical foundations of inventory management are mature, anchored by the optimality of $(s, S)$ policies for single-echelon systems with linear costs \citep{scarf1960optimality} and their efficient computation \citep{axsater2006inventory}. Interested readers may refer to textbooks such as \cite{simchi2005logic} for further details. However, real-world complexities—such as lead-time variability, lost sales, and multi-echelon networks—often render exact dynamic programming intractable due to the curse of dimensionality.
To address these high-dimensional state spaces, recent scholarship has pivoted toward Deep Reinforcement Learning (DRL). DRL, when combined with deep neural networks, effectively addresses the high-dimensional state and action spaces inherent in inventory control, alleviating the curse of dimensionality. Early studies demonstrated the feasibility of RL in inventory control. For instance, \cite{oroojlooyjadid2022deep} utilized the Deep Q-network (DQN) to solve the beer distribution game, a widely studied supply chain simulation. Similarly, \cite{gijsbrechts2022can} employed the A3C algorithm to achieve heuristic-level performance, while \cite{stranieri2023comparing} benchmarked multiple DRL methods, including A3C, PPO, and vanilla policy gradient (VPG), for inventory problems. Recent advancements have extended DRL to a variety of inventory control scenarios, such as managing non-stationary uncertain demand \citep{dehaybe2024deep, park2023adaptive}, optimizing multi-product systems \citep{sultana2020reinforcement, selukar2022inventory}, and handling diverse product types \citep{meisheri2020using, meisheri2022scalable}. DRL has also been applied to complex supply chain structures, including multi-echelon systems \citep{wu2023distributional, alvo2023neural, liu2024reinforcement, stranieri2024combining}, one-warehouse multi-retailer networks \citep{kaynov2024deep}, and the stochastic capacitated lot-sizing problem \citep{van2023using}.



Despite these algorithmic advances, a deployment gap remains. These methods require formal mathematical modeling and hyperparameter tuning that are inaccessible to the average small-business manager. Our work does not seek to invent a new DRL algorithm; rather, we leverage these existing powerful solvers (the ``Optimization Agent'') and focus on the \emph{interface} required to make them usable by non-experts.

\subsubsection{LLMs for Optimization and Decision Support.}
The emergence of LLMs has sparked intense interest in automating decision making (e.g. \cite{huang2025orlm}. LLMs are increasingly being adopted as versatile decision-support tools across many critical sectors. Within business and operations management, LLMs are leveraged to optimize processes, improve efficiency, and drive innovation \citep{li2024large}. Significant recent research interest has focused on their ability to automatically formulate optimization problems and dynamic programming problems from natural language descriptions \citep{ahmaditeshnizi2023optimus,ahmaditeshnizi2024optimus,zhou2025auto,liang2025llm}. In supply chain management specifically, LLMs are used for tasks such as demand forecasting and logistics optimization \citep{lu2024optimizing}, and recent work has also focused on developing agentic frameworks for advanced applications \citep{quan2024invagent,qi2025leveraging,long2025autonomous,simchi2026large,cohen2025supply}. In the healthcare domain, LLMs show remarkable potential to augment clinical decision-making by rapidly synthesizing patient information, generating differential diagnoses, and suggesting treatment plans, positioning them as powerful co-pilots for clinicians \citep{shah2023creation,nazi2024large}. Similarly, the financial sector uses LLMs for a wide range of applications, from market analysis to risk management, with specialized models like FinGPT and BloombergGPT enabling nuanced sentiment analysis, automated report generation, and enhanced algorithmic trading strategies.

Our work advances this direction by addressing the \emph{ambiguity} of the inputs. Existing frameworks typically assume the user provides a complete problem description. In contrast, our \emph{Information Extraction Agent} assumes the user is boundedly rational and the problem is initially ill-posed, requiring an iterative dialogue to elicit necessary parameters (e.g., distinguishing backorders from lost sales) before the solver can be invoked. This design makes our framework particularly suitable for small and medium-sized enterprises that lack in-house analytical staff but still face nontrivial inventory trade-offs. 




\subsubsection{Generative Agents and User Simulation.}
Evaluating interactive decision support systems presents a methodological dilemma: static datasets (e.g., standard NLP benchmarks) fail to capture the multi-turn dynamics of problem formulation, while human-subject trials are resource-intensive and difficult to reproduce. To address this, we draw upon the rich history of user simulation in dialogue systems and the recent emergence of generative agents.

User simulators have long been a staple in training task-oriented dialogue systems, particularly for Reinforcement Learning (RL) agents. Early approaches relied on agenda-based mechanisms \citep{schatzmann2007agenda}, where the simulator followed a strict stack of goals (e.g., ``book a flight'', ``specify time''). While effective for slot-filling tasks, these rule-based systems lacked the linguistic diversity of real users. The field subsequently moved toward data-driven approaches, utilizing sequence-to-sequence models to learn user behavior directly from corpora \citep{kreyssig2018neural, shi2019build}. However, these models often struggled with response collapse, generating generic or repetitive answers that failed to challenge the system.
The advent of Large Language Models has revolutionized user simulation by enabling ``Generative Agents'' that maintain consistent personas and memories. \cite{park2023generative} shows that LLMs could simulate believable social interactions in a sandbox environment. In the context of task-oriented dialogue, \cite{wang2025user} established that LLMs can act as zero-shot user simulators that outperform traditional models.
Crucially for our work, recent literature emphasizes simulating human imperfections. \cite{wang2024incharacter} investigated role-playing capabilities, finding that LLMs can effectively mimic specific demographic traits and knowledge gaps. \cite{li2023camel} introduced ``CAMEL'', a framework where two LLM agents (a user and an assistant) interact autonomously to solve tasks, revealing that agent-to-agent simulation can uncover edge cases that static testing misses.

Our \emph{Human Imitator} synthesizes these streams. Unlike the perfect ``agenda-based'' simulators of the past, our agent is designed to replicate the ``bounded rationality'' \citep{simon1955behavioral} of a non-expert supply chain manager. By conditioning the simulator on personas that exhibit specific knowledge gaps (e.g., confusing ``lost-sale'' with ``back-order''), we create a rigorous testbed that evaluates not just the system's ability to solve math, but its ability to clarify ambiguity. By fine-tuning a model on real human-machine interactions, we create a ``digital twin'' of a small business owner. This allows us to stress-test the system's ability to handle noise and ambiguity at scale, providing a more rigorous evaluation metric than simple code-generation accuracy.